% ============================================================
%  FRONT MATTER: Title, Authors, Abstract
% ============================================================

% ---- Title ----
\title{\LARGE\bfseries Merit Aid, Sorting, and Institutional Access:%
  \\[0.3em] Evidence from Tennessee's HOPE Scholarship%
  \thanks{We are grateful to
  \href{mailto:dg3342@columbia.edu}{Daniel Garcia},
  \href{mailto:fam2175@columbia.edu}{Francine Montecinos Puente}, \href{mailto:jackrosetti@gmail.com}{Jack Rosetti}, and
  \href{https://sites.google.com/view/seyhan-erden/home}{Seyhan Erden}
  (Columbia University);
  \href{https://www.hridaykarnani.com/}{Hriday Karnani}
  (Harvard University); and
  \href{mailto:as7596@columbia.edu}{Andrea Scalenghe}
  (MIT Sloan)
  for helpful discussions and comments. We also thank the Federal Reserve
  Bank of Cleveland and the Peer Review Board of the Economic Scholars
  Program. The views expressed herein are those of the authors and do not
  necessarily reflect the views of the Federal Reserve Bank of New York
  or Columbia University.}}

\author{%
  \vspace{-1.2em}\\
  {\large
  \href{mailto:bok2106@columbia.edu}{Balint Kidd} \quad
  \href{mailto:trs2169@columbia.edu}{Tommy Soltanian} \quad
  \href{mailto:tyler.sotomayor@columbia.edu}{Tyler Sotomayor} \quad
  \href{mailto:gu2145@columbia.edu}{Gabriel Uceda-Sosa}} \\[0.7em]
  {\large Columbia University} \\[0.7em]
  {\normalsize April 2026} \\[1.2em]
  {\small\color{blue}\href{https://raw.githubusercontent.com/tylersotomayor/merit-aid-working-paper/main/main.pdf}{Link to Most Current Version}}
}

\date{\vspace{-3em}}

% ============================================================
\begin{document}
% ============================================================

\workingpaperheader

\begingroup
\onehalfspacing
\maketitle

\vspace{1em}

\begin{abstract}
\noindent
Merit-based scholarship programs are widely viewed as regressive
because eligibility correlates with family income. Using tax-linked
data from \textcite{chetty2020mobility}, we show that Tennessee's HOPE
Scholarship produces a \emph{progressive} compositional shift at
public colleges: all five income quintiles shift monotonically, with
bottom-quintile share rising and top-quintile share falling. Total
enrollment is unchanged---the shift reflects sorting, not expansion.
We decompose this effect into a permanent substitution effect (the
scholarship's price reduction) and a cyclical income effect
(sensitivity to economic downturns). Merging county-level income
data, we find that tuition-scaling programs (Florida, West Virginia)
compensate the income effect---composition is stable during
downturns---while fixed-dollar programs (Tennessee, Georgia) do not,
as marginal students reallocate across tiers. This produces a second
monotonic quintile gradient estimated on entirely different variation.
Cross-state evidence confirms the model: Massachusetts (progressive,
as predicted), South Dakota (regressive, as predicted). Program
design---eligibility threshold, award structure, and need-based aid
interaction---jointly predicts both the direction and durability of
compositional change.
\end{abstract}

\vspace{1em}

\noindent \textbf{JEL Codes:} I23, I24, I28, H75, J24, C23

\vspace{0.25em}

\noindent \textbf{Keywords:} merit aid, college composition, sorting, HOPE Scholarship, difference-in-differences

\endgroup

\doublespacing
\newpage

\tableofcontents
\newpage
